\scanpage{34}
whose kernel \inserted{is $\bigcap_n nH^0(K,G)$, hence} consists of the elements of $G(K)$ \uline{infinitely divisible} in $G$.

Now here $K$ being an \deleted{group} \inserted{absolutely} finitely generated field, the Mordell-Weil theorem assures us that $G(K)$ is \deleted{a group} \inserted{a $\mathbb Z$-module of} finite type -- hence $G(K)\longrightarrow\varprojlim G(K)_n$ is injective. Thus $G(K)\longrightarrow\varprojlim H^1(K,{}_nG)$ is injective, and a fortiori the homomorphism (16), which was to be proved.

\uline{Remark} It would not suffice in this argument to consider, for one $\ell$, the extension of $E_{\overline K/K}$ by $H_1(X,\mathbb Z_\ell)$ -- one must use all the $\ell$ -- if one omits $p$, then translation by a point of order $p$ gives a counterexample. \deleted{\unsure{Thus}} The result amounts to this: For every $x\in G(K)$, we consider the torsor over $K$, with group $T_\infty(G)=\varprojlim {}_nG$, image
\margin{33}
